\documentclass[10pt,a4paper]{article}

% ==============================================================================
% TEMPLATE EDITORIAL — BOLETIM FLASH (JORNAL ECONÔMICO DO TOCANTINS)
% Estilo: Briefing Executivo de Jornal Econômico (2 Páginas Balanceadas)
% Paleta FLASH: Azul-Marinho (#0F2537), Azul Secundário (#244B6B),
%               Dourado Institucional (#C8942B) e Fundo Neutro (#F4F6F9)
% ==============================================================================
\usepackage[utf8]{inputenc}
\usepackage[T1]{fontenc}
\usepackage[brazil]{babel}
\usepackage{helvet}
\renewcommand{\familydefault}{\sfdefault}

\usepackage[top=1.35cm, bottom=1.45cm, left=1.40cm, right=1.40cm, headheight=14pt, headsep=8pt, footskip=18pt]{geometry}
\usepackage{xcolor}
\usepackage{graphicx}
\usepackage{booktabs}
\usepackage{tabularx}
\usepackage{longtable}
\usepackage{array}
\usepackage{colortbl}
\usepackage{tikz}
\usepackage{fancyhdr}
\usepackage{lastpage}
\usepackage{enumitem}
\usepackage{microtype}
\usepackage{setspace}
\usepackage{multicol}
\usepackage{amssymb}
\usepackage{hyperref}

% Cores institucionais do Boletim FLASH
\definecolor{corprimaria}{HTML}{0F2537}
\definecolor{corsecundaria}{HTML}{244B6B}
\definecolor{corfundo}{HTML}{F4F6F9}
\definecolor{douradoto}{HTML}{C8942B}
\definecolor{coracento}{HTML}{C8942B}
\definecolor{cinzatexto}{HTML}{1E293B}
\definecolor{cinzasub}{HTML}{475569}
\definecolor{cinzaborda}{HTML}{CBD5E1}
\definecolor{verdealta}{HTML}{15803D}
\definecolor{vermelhobaixa}{HTML}{B91C1C}

\hypersetup{
  colorlinks=true,
  linkcolor=corprimaria,
  urlcolor=corsecundaria,
  pdftitle={Boletim Flash — PIB Trimestral do Tocantins e Comparativo Setorial com o Brasil},
  pdfauthor={Jornal Economico do Tocantins}
}

\setlength{\parindent}{0pt}
\setlength{\parskip}{3.5pt}
\setlength{\columnsep}{0.55cm}
\setlength{\columnseprule}{0.4pt}
\def\columnseprulecolor{\color{cinzaborda}}
\renewcommand{\arraystretch}{1.15}

% Cabeçalho e Rodapé de Jornal Econômico (Sem colisão central)
\pagestyle{fancy}
\fancyhf{}
\renewcommand{\headrulewidth}{0pt}
\renewcommand{\footrulewidth}{0.5pt}
\renewcommand{\footrule}{\hbox to\headwidth{\color{douradoto}\leaders\hrule height \footrulewidth\hfill}}

\fancyhead[L]{\ifnum\value{page}>1 \scriptsize\bfseries\color{corprimaria}JORNAL ECONÔMICO DO TOCANTINS \enspace{\color{douradoto}$\vert$}\enspace {\mdseries\color{cinzasub}BOLETIM RÁPIDO (FLASH MENSAL) $\cdot$ 1\textordmasculine{} Trimestre de 2026}\fi}
\fancyhead[R]{\ifnum\value{page}>1 \scriptsize\color{cinzasub}\#12/2026\fi}

\fancyfoot[L]{\scriptsize\color{cinzasub}\textbf{\color{corprimaria}Jornal Econômico do Tocantins} $\cdot$ Produto Interno Bruto e Contas Regionais do Tocantins $\cdot$ Publicado em 14/11/2026}
\fancyfoot[R]{\scriptsize\bfseries\color{corprimaria}Pág. \thepage\ de \pageref*{LastPage}}

% Comando para Títulos de Seção Estilo Jornal Financeiro
\newcommand{\secaoboletim}[1]{%
  \par\vspace{0.16cm}%
  \noindent\begin{tikzpicture}
    \fill[corprimaria] (0,0) rectangle (\textwidth, 0.52);
    \fill[douradoto] (0,0) rectangle (0.16, 0.52);
    \node[anchor=west, font=\small\bfseries, text=white] at (0.30, 0.26) {\MakeUppercase{#1}};
  \end{tikzpicture}%
  \par\vspace{0.08cm}%
}

\newcommand{\subsecaoboletim}[1]{%
  \par\vspace{0.10cm}%
  \noindent{\small\bfseries\color{corprimaria}\rule{0.14cm}{0.26cm}\hspace{5pt}#1}%
  \par\vspace{0.03cm}%
  {\color{douradoto}\hrule height 0.6pt}%
  \vspace{0.06cm}%
}

\begin{document}
\color{cinzatexto}

% ==============================================================================
% MASTHEAD DE JORNAL ECONÔMICO (AUTO-AJUSTÁVEL, SEM SOBREPOSIÇÃO DE TEXTO)
% ==============================================================================
\noindent\begin{tikzpicture}
  \node[
    fill=corprimaria,
    inner xsep=12pt,
    inner ysep=10pt,
    text width=\textwidth-24pt,
    align=left
  ] (masthead) {
    % Linha superior do cabeçalho do jornal
    \hbox to \linewidth{%
      {\colorbox{douradoto}{\strut\hspace{4pt}\scriptsize\bfseries\color{corprimaria}BOLETIM RÁPIDO (FLASH MENSAL) $\cdot$ \#12/2026\hspace{4pt}}}%
      \hfill
      {\scriptsize\color{white}\textbf{JORNAL ECONÔMICO DO TOCANTINS} \enspace{\color{douradoto}$\vert$}\enspace Ref.: \textbf{1\textordmasculine{} Trimestre de 2026} \enspace{\color{douradoto}$\vert$}\enspace 14/11/2026}%
    }%
    \par\vspace{5pt}%
    {\color{douradoto}\hrule height 0.8pt}%
    \par\vspace{6pt}%
    % Título Principal e Subtítulo com quebra automática limpa
    {\fontsize{13.5}{16.2}\selectfont\bfseries\color{white}Boletim Flash — PIB Trimestral do Tocantins e Comparativo Setorial com o Brasil\par}%
    \vspace{3.5pt}%
    {\footnotesize\color{corfundo}Pareamento lado a lado entre a série dessazonalizada X-13ARIMA-SEATS (SEPLAN-TO) e a Tabela 1621 do IBGE, acrescido das Contas Regionais (Tabela 5938)\par}%
  };
  \fill[douradoto] (masthead.south west) rectangle ([yshift=-2.5pt]masthead.south east);
\end{tikzpicture}

\vspace{0.14cm}

% ==============================================================================
% 1. FAIXA DE INDICADORES-CHAVE DO PAINEL (4 CARDS KPI EM LINHA ÚNICA)
% ==============================================================================
\noindent\makebox[\textwidth][s]{%
\begin{tikzpicture}[baseline=(box.base)]
  \node[fill=corfundo, draw=cinzaborda, line width=0.45pt, rounded corners=2pt,
        inner xsep=6pt, inner ysep=4pt, text width=0.286\textwidth, minimum height=1.22cm, align=left] (box) {%
    {\fontsize{6.2}{7.5}\selectfont\bfseries\color{corsecundaria}PIB NOMINAL (TOCANTINS)}\par\vspace{1.2pt}%
    {\fontsize{9.0}{10.8}\selectfont\bfseries\color{corprimaria}R\$ 25,8 Bi $\blacktriangledown$ 1,3\%}\par\vspace{1pt}%
    {\fontsize{5.9}{7.0}\selectfont\color{cinzasub}Ref.: T1 2026 (SEPLAN-TO)}%
  };
  \fill[coracento] (box.north west) rectangle ([xshift=2.4pt]box.south west);
\end{tikzpicture}\hfill%
\begin{tikzpicture}[baseline=(box.base)]
  \node[fill=corfundo, draw=cinzaborda, line width=0.45pt, rounded corners=2pt,
        inner xsep=6pt, inner ysep=4pt, text width=0.286\textwidth, minimum height=1.22cm, align=left] (box) {%
    {\fontsize{6.2}{7.5}\selectfont\bfseries\color{corsecundaria}IMPOSTOS LÍQUIDOS (TO)}\par\vspace{1.2pt}%
    {\fontsize{9.0}{10.8}\selectfont\bfseries\color{corprimaria}R\$ 2,26 Bi $\blacktriangledown$ 0,2\%}\par\vspace{1pt}%
    {\fontsize{5.9}{7.0}\selectfont\color{cinzasub}Preços Correntes e Var. Dessaz.}%
  };
  \fill[coracento] (box.north west) rectangle ([xshift=2.4pt]box.south west);
\end{tikzpicture}\hfill%
\begin{tikzpicture}[baseline=(box.base)]
  \node[fill=corfundo, draw=cinzaborda, line width=0.45pt, rounded corners=2pt,
        inner xsep=6pt, inner ysep=4pt, text width=0.286\textwidth, minimum height=1.22cm, align=left] (box) {%
    {\fontsize{6.2}{7.5}\selectfont\bfseries\color{corsecundaria}PIB NOMINAL (BRASIL)}\par\vspace{1.2pt}%
    {\fontsize{9.0}{10.8}\selectfont\bfseries\color{corprimaria}R\$ 3,43 Tri $\blacktriangle$ 0,5\%}\par\vspace{1pt}%
    {\fontsize{5.9}{7.0}\selectfont\color{cinzasub}IBGE Tabelas 1846 / 5932}%
  };
  \fill[coracento] (box.north west) rectangle ([xshift=2.4pt]box.south west);
\end{tikzpicture}%
}
\par\vspace{0.03cm}
\noindent{\fontsize{6.3}{7.5}\selectfont\bfseries\color{corsecundaria}DESEMPENHO SETORIAL NO TRIMESTRE (DO MAIOR CRESCIMENTO AO MENOR, DA ESQUERDA PARA A DIREITA)}\par\vspace{0.02cm}
\noindent\makebox[\textwidth][s]{%
\begin{tikzpicture}[baseline=(box.base)]
  \node[fill=corfundo, draw=cinzaborda, line width=0.45pt, rounded corners=2pt,
        inner xsep=6pt, inner ysep=4pt, text width=0.286\textwidth, minimum height=0.92cm, align=left] (box) {%
    {\fontsize{6.2}{7.5}\selectfont\bfseries\color{corsecundaria}1\textordmasculine{} Indústria}\par\vspace{1.2pt}%
    {\fontsize{9.0}{10.8}\selectfont\bfseries\color{verdealta}$\blacktriangle$ + 3,1\%}\par\vspace{1pt}%
    {\fontsize{5.9}{7.0}\selectfont\color{cinzasub}Maior crescimento trimestral}%
  };
  \fill[verdealta] (box.north west) rectangle ([xshift=2.4pt]box.south west);
\end{tikzpicture}\hfill%
\begin{tikzpicture}[baseline=(box.base)]
  \node[fill=corfundo, draw=cinzaborda, line width=0.45pt, rounded corners=2pt,
        inner xsep=6pt, inner ysep=4pt, text width=0.286\textwidth, minimum height=0.92cm, align=left] (box) {%
    {\fontsize{6.2}{7.5}\selectfont\bfseries\color{corsecundaria}2\textordmasculine{} Serviços}\par\vspace{1.2pt}%
    {\fontsize{9.0}{10.8}\selectfont\bfseries\color{verdealta}$\blacktriangle$ + 1,5\%}\par\vspace{1pt}%
    {\fontsize{5.9}{7.0}\selectfont\color{cinzasub}Desempenho positivo na margem}%
  };
  \fill[verdealta] (box.north west) rectangle ([xshift=2.4pt]box.south west);
\end{tikzpicture}\hfill%
\begin{tikzpicture}[baseline=(box.base)]
  \node[fill=corfundo, draw=cinzaborda, line width=0.45pt, rounded corners=2pt,
        inner xsep=6pt, inner ysep=4pt, text width=0.286\textwidth, minimum height=0.92cm, align=left] (box) {%
    {\fontsize{6.2}{7.5}\selectfont\bfseries\color{corsecundaria}3\textordmasculine{} Agropecuária}\par\vspace{1.2pt}%
    {\fontsize{9.0}{10.8}\selectfont\bfseries\color{vermelhobaixa}$\blacktriangledown$ - 7,2\%}\par\vspace{1pt}%
    {\fontsize{5.9}{7.0}\selectfont\color{cinzasub}Retração trimestral na margem}%
  };
  \fill[vermelhobaixa] (box.north west) rectangle ([xshift=2.4pt]box.south west);
\end{tikzpicture}%
}


\vspace{0.10cm}

% ==============================================================================
% 2. ANÁLISE PRELIMINAR AUTOMATIZADA DO SCRIPT (RASTREAMENTO DE VARIAÇÕES)
% ==============================================================================
\noindent\begin{tikzpicture}
  \node[fill=corfundo, draw=cinzaborda, line width=0.5pt, rounded corners=2pt,
        inner xsep=8pt, inner ysep=5pt, text width=\textwidth-17pt, font=\scriptsize\color{cinzatexto}] (dbox) {
    {\footnotesize\bfseries\color{corprimaria}Síntese de Variações entre Períodos Subsequentes}\par\vspace{2pt}
    \begin{itemize}[leftmargin=11pt, itemsep=1.0pt, topsep=1pt, parsep=0pt]
       \item \textbf{\color{vermelhobaixa}[OSCILAÇÃO NEGATIVA LEVE] PIB Trimestral do Tocantins (Ajuste Sazonal X-13):} PIB Trimestral do Tocantins (Ajuste Sazonal X-13) apresentou acomodação negativa de -1,28\%. [nível: R\$ 25,8 Bi $\blacktriangledown$ 1,3\%]
       \item \textbf{\color{verdealta}[AVANÇO MODERADO] PIB Trimestral do Brasil (IBGE Tabela 1621):} PIB Trimestral do Brasil (IBGE Tabela 1621) avançou +1,07\% em trajetória de expansão gradual. [nível: R\$ 3,43 Tri $\blacktriangle$ 0,5\%]
       \item \textbf{\color{verdealta}[AVANÇO MODERADO] 1\textordmasculine{} Setor em Desempenho (Indústria):} 1\textordmasculine{} Setor em Desempenho (Indústria) avançou +3,07\% em trajetória de expansão gradual. [nível: $\blacktriangle$ + 3,1\%]
       \item \textbf{\color{vermelhobaixa}[RETRAÇÃO RELEVANTE] 3\textordmasculine{} Setor em Desempenho (Agropecuária):} 3\textordmasculine{} Setor em Desempenho (Agropecuária) recuou -7,16\%, sinalizando desaceleração relevante no período. [nível: $\blacktriangledown$ - 7,2\%]
    \end{itemize}
  };
  \fill[corprimaria] (dbox.north west) rectangle ([xshift=3pt]dbox.south west);
\end{tikzpicture}


\vspace{0.08cm}

% ==============================================================================
% 3. ANÁLISE CONJUNTURAL DA EQUIPE E PERSPECTIVAS (2 COLUNAS DE JORNAL)
% ==============================================================================
\secaoboletim{1. Análise Conjuntural da Equipe e Cenários}

\begin{multicols}{2}
\begin{spacing}{1.06}
{\small A comparação direta entre a série trimestral do Tocantins — ajustada sazonalmente pelo método X-13ARIMA-SEATS — e a Tabela 1621 do IBGE para o Brasil evidencia uma mudança importante nos vetores de tração de curto prazo da economia estadual. No trimestre mais recente contra o trimestre imediatamente anterior, a Indústria tocantinense avançou 3,1\% e o setor de Serviços cresceu 1,5\%, superando o ritmo médio nacional em ambos os grupamentos.

\medskip
Por sua vez, as Contas Regionais do Brasil (Tabela 5938 do IBGE) confirmam o ganho estrutural de participação do Tocantins tanto na Região Norte quanto no Produto Interno Bruto brasileiro, sustentado pela expansão contínua do Valor Adicionado Bruto da agropecuária, da agroindústria e dos serviços logísticos.}

\subsecaoboletim{Perspectivas e Cenários}
{\small O desempenho positivo dos setores secundário e terciário na margem sinaliza maior diversificação produtiva intraestadual, reduzindo a volatilidade macroeconômica nos trimestres de entressafra agrícola.}

\end{spacing}
\end{multicols}

% ==============================================================================
% 4. PANORAMA ESTATÍSTICO, GRÁFICO E TABELAS EXECUTIVAS DO PERÍODO
% ==============================================================================
\secaoboletim{2. Panorama Estatístico, Gráfico e Tabelas Executivas do Período}

\begin{center}\includegraphics[width=0.95\textwidth]{figuras/2026_11_flash_pib_pib_to_br.pdf}\end{center}
\par\vspace{0.08cm}
\noindent\begin{minipage}{\linewidth}
  \noindent{\fontsize{7.8}{9.5}\selectfont\bfseries\color{corprimaria}\rule{0.10cm}{0.20cm}\hspace{4pt}Comparativo Lado a Lado: Série Dessazonalizada X-13 (Tocantins) vs. Tabela 1621 IBGE (Brasil)}\par\vspace{2pt}
  {\scriptsize\setlength{\tabcolsep}{4.5pt}%
\noindent\resizebox{\linewidth}{!}{%
\begin{tabular}{lrrrrr}
\toprule
\rowcolor{corprimaria} \textbf{\color{white}Trimestre} & \textbf{\color{white}PIB TO: Índ. (Var.)} & \textbf{\color{white}PIB BR: Índ. (Var.)} & \textbf{\color{white}Agro TO/BR} & \textbf{\color{white}Ind. TO/BR} & \textbf{\color{white}Serv. TO/BR} \\
\midrule
T4 2024 & 177.2 (+1.64\%) & 191.4 (+0.08\%) & +17.1\% / -2.2\% & +0.1\% / +0.3\% & -0.7\% / +0.0\% \\
\rowcolor{corfundo} T1 2025 & 183.5 (+3.54\%) & 193.8 (+1.26\%) & +3.5\% / +16.1\% & -3.3\% / +0.0\% & +2.1\% / +0.2\% \\
T2 2025 & 180.1 (-1.85\%) & 194.5 (+0.36\%) & -1.1\% / -3.2\% & +5.9\% / +0.7\% & -1.1\% / +0.8\% \\
\rowcolor{corfundo} T3 2025 & 179.7 (-0.20\%) & 194.7 (+0.11\%) & -4.6\% / +0.5\% & -0.7\% / +0.6\% & +0.5\% / +0.3\% \\
T4 2025 & 181.6 (+1.01\%) & 195.2 (+0.25\%) & +6.5\% / +0.7\% & -3.3\% / -0.7\% & +1.4\% / +0.6\% \\
\rowcolor{corfundo} T1 2026 & 179.2 (-1.28\%) & 197.3 (+1.07\%) & -7.2\% / +2.2\% & +3.1\% / +1.2\% & +1.5\% / +0.4\% \\
\bottomrule
\end{tabular}%
}%
}
  \par\vspace{1.5pt}{\fontsize{6.0}{7.2}\selectfont\color{cinzasub}\textit{Fonte: SEPLAN-TO (ajuste sazonal próprio X-13ARIMA-SEATS) e IBGE/SIDRA (Tabela 1621: códs. 90687, 90691, 90696 e 90707).}}
\end{minipage}
\par\vspace{0.08cm}

\par\vspace{0.08cm}
\noindent\begin{minipage}{\linewidth}
  \noindent{\fontsize{7.8}{9.5}\selectfont\bfseries\color{corprimaria}\rule{0.10cm}{0.20cm}\hspace{4pt}Contas Regionais do Brasil (IBGE - Tabela 5938): Estrutura do PIB e Valor Adicionado do Tocantins}\par\vspace{2pt}
  {\scriptsize\setlength{\tabcolsep}{4.5pt}%
\noindent\resizebox{\linewidth}{!}{%
\begin{tabular}{lrrrrrr}
\toprule
\rowcolor{corprimaria} \textbf{\color{white}Ano} & \textbf{\color{white}PIB (R\$ Mi)} & \textbf{\color{white}Impostos (R\$ Mi)} & \textbf{\color{white}VA Agro (R\$ Mi)} & \textbf{\color{white}VA Ind. (R\$ Mi)} & \textbf{\color{white}VA Serv. (R\$ Mi)} & \textbf{\color{white}Part. Norte / BR} \\
\midrule
2019 & R\$ 39.355,9 & R\$ 3.817,5 & R\$ 5.036,9 & R\$ 4.211,5 & R\$ 26.290,1 & 9.36\% / 0.53\% \\
\rowcolor{corfundo} 2020 & R\$ 43.649,8 & R\$ 4.140,9 & R\$ 8.033,2 & R\$ 4.408,4 & R\$ 27.067,3 & 9.13\% / 0.57\% \\
2021 & R\$ 51.780,8 & R\$ 5.085,4 & R\$ 13.364,4 & R\$ 5.000,3 & R\$ 28.330,7 & 9.18\% / 0.57\% \\
\rowcolor{corfundo} 2022 & R\$ 58.208,8 & R\$ - & R\$ - & R\$ - & R\$ - & 10.13\% / 0.58\% \\
2023 & R\$ 64.317,7 & R\$ - & R\$ - & R\$ - & R\$ - & 10.10\% / 0.59\% \\
\bottomrule
\end{tabular}%
}%
}
  \par\vspace{1.5pt}{\fontsize{6.0}{7.2}\selectfont\color{cinzasub}\textit{Fonte: IBGE — Contas Regionais do Brasil (Tabela 5938). Nota: Pesquisa estrutural anual com defasagem oficial de \textasciitilde{}2 anos.}}
\end{minipage}
\par\vspace{0.08cm}



\vfill

% ==============================================================================
% EXPEDIENTE EDITORIAL E NOTA METODOLÓGICA
% ==============================================================================
\noindent\begin{tikzpicture}
  \node[
    fill=corfundo,
    draw=cinzaborda,
    line width=0.5pt,
    rounded corners=2pt,
    inner sep=6pt,
    text width=\textwidth-13pt,
    font=\scriptsize\color{cinzasub}
  ] {
    \textbf{\color{corprimaria}Nota Metodológica e Fontes Primárias:} Séries trimestrais do Tocantins dessazonalizadas pelo algoritmo oficial X-13ARIMA-SEATS do U.S. Census Bureau, comparadas lado a lado com a Tabela 1621 do IBGE (Contas Nacionais Trimestrais) e complementadas pelas Contas Regionais do Brasil (Tabela 5938 do IBGE).\par\vspace{2pt}
    \textbf{\color{corprimaria}Expediente Editorial:} Lorenzo Costa | Graduado em Economia — Jornal Econômico do Tocantins \enspace$\vert$\enspace \textbf{\color{corprimaria}Publicação:} Núcleo de Conjuntura Econômica — \textit{Jornal Econômico do Tocantins}.
  };
\end{tikzpicture}

\end{document}
